%% Согласно ГОСТ Р 7.0.11-2011:
%% 5.3.3 В заключении диссертации излагают итоги выполненного исследования, рекомендации, перспективы дальнейшей разработки темы.
%% 9.2.3 В заключении автореферата диссертации излагают итоги данного исследования, рекомендации и перспективы дальнейшей разработки темы.
\begin{enumerate}
    \item Задача декодирования визуальной информации по сигналам мозга формализована в единой постановке, объединяющей три уровня описания структуры данных: временную динамику, пространственно-временную структуру и мультимодальное объединение.
    \item Предложена линейная авторегрессионная модель прогнозирования фМРТ-отклика по визуальному стимулу; полученная оценка гемодинамического лага согласуется с известным физиологическим диапазоном.
    \item Разработан метод пространственно-временной классификации временных рядов фМРТ на основе масок активности и римановой геометрии ковариационных матриц, статистически значимо превосходящий нейросетевые аналоги в условиях ограниченной выборки.
    \item Предложена мультимодальная архитектура реконструкции визуальных стимулов по одновременным сигналам фМРТ и ЭЭГ; по метрике CLIP-Score она превосходит одномодальные аналоги, а априорная диффузионная модель дополнительно повышает семантическую согласованность реконструкций.
    \item Кросс-доменные эксперименты очертили границы применимости предложенных методов и обозначили доменную адаптацию как направление дальнейшей работы.
\end{enumerate}
